% GENERATED FILE. Edit canonical lessons, then run npm run generate:books.
% canonical-source-hash: fnv1a64:1b3fda6b398a90f5
% canonical-lessons: AR-C387-hijra, AR-C387-umma, AR-C387-jinsiyya, AR-C387-aghlabiyya, AR-C387-aqalliyya

\chapter{People and nations}
\label{ch:ar-people-and-nations}
\hlchaptermodality{387}

\begin{chapteropening}
\textbf{By the end of this chapter:} \emph{I can talk about migration, a nation, nationality, a majority and a minority.}

Complete the last lesson of chapter 387: I can talk about migration, a nation, nationality, a majority and a minority.
\end{chapteropening}

\section[\texorpdfstring{\ar{هجرة} (hijra)}{هجرة (hijra)}]{\ar{هجرة} (hijra) — migration}

\label{lesson:AR-C387-hijra}



\begin{quote}
Before the new one: say the Arabic for a conflict, then the Arabic for a crisis.
\end{quote}

\subsection*{You'll want to know: \ar{هجرة}}
\textbf{\ar{هجرة}} — \emph{hijra} — \textquotedblleft{}migration\textquotedblright{}.

\begin{quote}
\ar{زادت} \ar{الهجرة} \ar{إلى} \ar{المدن}.
\end{quote}

\begin{quote}
\emph{zādat al-hijra ilā al-mudun.}
\end{quote}

\begin{quote}
Migration to the cities increased.
\end{quote}

\begin{cousinweb}
\textbf{\ar{هجرة}} (\emph{hijra}) is built on the native Arabic root \emph{h-j-r}, \textquotedblleft{}to leave, to cut off\textquotedblright{} (Lane). \emph{Hijra} is leaving one's land for another: migration.
\end{cousinweb}

\begin{grammarlens}[title={What you've built}]
One more word for giving reasons and opinions.
\end{grammarlens}

\subsection*{Your turn}
\begin{itemize}
\raggedright
  \item \emph{Say it:} \emph{hijra}
  \item \emph{Say it:} \emph{hijra}, once more
  \item \emph{Say it:} \emph{zādat al-hijra ilā al-mudun.} — Migration to the cities increased.
  \item \emph{Recall:} say the Arabic for a conflict, then the Arabic for a crisis, then say \emph{hijra} again
\end{itemize}

\begin{tcolorbox}[breakable,colback=teal!4,colframe=teal!35!black,title={Before you move on}]
What does \ar{هجرة} mean? (\textquotedblleft{}Migration\textquotedblright{}.) Say it once more.
\end{tcolorbox}
\section[\texorpdfstring{\ar{أمة} (umma)}{أمة (umma)}]{\ar{أمة} (umma) — a nation}

\label{lesson:AR-C387-umma}



\begin{quote}
Before the new one: say the Arabic for a crisis, then the Arabic for migration.
\end{quote}

\subsection*{You'll want to know: \ar{أمة}}
\textbf{\ar{أمة}} — \emph{umma} — \textquotedblleft{}a nation\textquotedblright{}.

\begin{quote}
\ar{تحتفل} \ar{الأمة} \ar{بعيدها} \ar{الوطني}.
\end{quote}

\begin{quote}
\emph{taḥtafilu al-umma bi-ʿīdihā al-waṭanī.}
\end{quote}

\begin{quote}
The nation is celebrating its national day.
\end{quote}

Plural: \ar{أمم}, \emph{umam}.

\begin{cousinweb}
\textbf{\ar{أمة}} (\emph{umma}) is built on the native Arabic root \emph{ʾ-m-m}, \textquotedblleft{}to aim at, to head for\textquotedblright{} (Lane). An \emph{umma} is a people that head the same way: a nation. \ar{أم}, mother, is from the same root.
\end{cousinweb}

\begin{grammarlens}[title={What you've built}]
One more word for giving reasons and opinions.
\end{grammarlens}

\subsection*{Your turn}
\begin{itemize}
\raggedright
  \item \emph{Say it:} \emph{umma}
  \item \emph{Say it:} \emph{umma}, once more
  \item \emph{Say it:} \emph{taḥtafilu al-umma bi-ʿīdihā al-waṭanī.} — The nation is celebrating its national day.
  \item \emph{Recall:} say the Arabic for a crisis, then the Arabic for migration, then say \emph{umma} again
\end{itemize}

\begin{tcolorbox}[breakable,colback=teal!4,colframe=teal!35!black,title={Before you move on}]
What does \ar{أمة} mean? (\textquotedblleft{}A nation\textquotedblright{}.) Say it once more.
\end{tcolorbox}
\section[\texorpdfstring{\ar{جنسية} (jinsiyya)}{جنسية (jinsiyya)}]{\ar{جنسية} (jinsiyya) — nationality}

\label{lesson:AR-C387-jinsiyya}



\begin{quote}
Before the new one: say the Arabic for migration, then the Arabic for a nation.
\end{quote}

\subsection*{You'll want to know: \ar{جنسية}}
\textbf{\ar{جنسية}} — \emph{jinsiyya} — \textquotedblleft{}nationality\textquotedblright{}.

\begin{quote}
\ar{يحمل} \ar{صديقي} \ar{جنسيتين}.
\end{quote}

\begin{quote}
\emph{yaḥmilu ṣadīqī jinsiyyatayn.}
\end{quote}

\begin{quote}
My friend holds two nationalities.
\end{quote}

\begin{cousinweb}
\textbf{\ar{جنسية}} (\emph{jinsiyya}) is built on the native Arabic root \emph{j-n-s}, \textquotedblleft{}a kind, a sort\textquotedblright{} (Lane). \emph{Jinsiyya} is a modern abstract noun from \emph{jins}, kind: the kind of citizen one is, nationality.
\end{cousinweb}

\begin{grammarlens}[title={What you've built}]
One more word for giving reasons and opinions.
\end{grammarlens}

\subsection*{Your turn}
\begin{itemize}
\raggedright
  \item \emph{Say it:} \emph{jinsiyya}
  \item \emph{Say it:} \emph{jinsiyya}, once more
  \item \emph{Say it:} \emph{yaḥmilu ṣadīqī jinsiyyatayn.} — My friend holds two nationalities.
  \item \emph{Recall:} say the Arabic for migration, then the Arabic for a nation, then say \emph{jinsiyya} again
\end{itemize}

\begin{tcolorbox}[breakable,colback=teal!4,colframe=teal!35!black,title={Before you move on}]
What does \ar{جنسية} mean? (\textquotedblleft{}Nationality\textquotedblright{}.) Say it once more.
\end{tcolorbox}
\section[\texorpdfstring{\ar{أغلبية} (aghlabiyya)}{أغلبية (aghlabiyya)}]{\ar{أغلبية} (aghlabiyya) — a majority}

\label{lesson:AR-C387-aghlabiyya}



\begin{quote}
Before the new one: say the Arabic for a nation, then the Arabic for nationality.
\end{quote}

\subsection*{You'll want to know: \ar{أغلبية}}
\textbf{\ar{أغلبية}} — \emph{aghlabiyya} — \textquotedblleft{}a majority\textquotedblright{}.

\begin{quote}
\ar{صوتت} \ar{الأغلبية} \ar{ضد} \ar{المشروع}.
\end{quote}

\begin{quote}
\emph{ṣawwatat al-aghlabiyya ḍidda al-mashrūʿ.}
\end{quote}

\begin{quote}
The majority voted against the project.
\end{quote}

\begin{cousinweb}
\textbf{\ar{أغلبية}} (\emph{aghlabiyya}) is built on the native Arabic root \emph{gh-l-b}, \textquotedblleft{}to overcome, to prevail\textquotedblright{} (Lane). \emph{Aghlabiyya} is a modern abstract noun from \emph{aghlab}, most: the greater part, a majority.
\end{cousinweb}

\begin{grammarlens}[title={What you've built}]
One more word for giving reasons and opinions.
\end{grammarlens}

\subsection*{Your turn}
\begin{itemize}
\raggedright
  \item \emph{Say it:} \emph{aghlabiyya}
  \item \emph{Say it:} \emph{aghlabiyya}, once more
  \item \emph{Say it:} \emph{ṣawwatat al-aghlabiyya ḍidda al-mashrūʿ.} — The majority voted against the project.
  \item \emph{Recall:} say the Arabic for a nation, then the Arabic for nationality, then say \emph{aghlabiyya} again
\end{itemize}

\begin{tcolorbox}[breakable,colback=teal!4,colframe=teal!35!black,title={Before you move on}]
What does \ar{أغلبية} mean? (\textquotedblleft{}A majority\textquotedblright{}.) Say it once more.
\end{tcolorbox}
\section[\texorpdfstring{\ar{أقلية} (aqalliyya)}{أقلية (aqalliyya)}]{\ar{أقلية} (aqalliyya) — a minority}

\label{lesson:AR-C387-aqalliyya}



\begin{quote}
Before the new one: say the Arabic for nationality, then the Arabic for a majority.
\end{quote}

\subsection*{You'll want to know: \ar{أقلية}}
\textbf{\ar{أقلية}} — \emph{aqalliyya} — \textquotedblleft{}a minority\textquotedblright{}.

\begin{quote}
\ar{تعيش} \ar{في} \ar{البلد} \ar{أقلية} \ar{صغيرة}.
\end{quote}

\begin{quote}
\emph{taʿīshu fī al-balad aqalliyya ṣaghīra.}
\end{quote}

\begin{quote}
A small minority lives in the country.
\end{quote}

\begin{cousinweb}
\textbf{\ar{أقلية}} (\emph{aqalliyya}) is built on the native Arabic root \emph{q-l-l}, \textquotedblleft{}to be few, little\textquotedblright{} (Lane). \emph{Aqalliyya} is a modern abstract noun from \emph{aqall}, fewer: a minority.
\end{cousinweb}

\begin{grammarlens}[title={What you've built}]
One more word for giving reasons and opinions.
\end{grammarlens}

\subsection*{Your turn}
\begin{itemize}
\raggedright
  \item \emph{Say it:} \emph{aqalliyya}
  \item \emph{Say it:} \emph{aqalliyya}, once more
  \item \emph{Say it:} \emph{taʿīshu fī al-balad aqalliyya ṣaghīra.} — A small minority lives in the country.
  \item \emph{Recall:} say the Arabic for nationality, then the Arabic for a majority, then say \emph{aqalliyya} again
\end{itemize}

\begin{tcolorbox}[breakable,colback=teal!4,colframe=teal!35!black,title={Before you move on}]
What does \ar{أقلية} mean? (\textquotedblleft{}A minority\textquotedblright{}.) Say it once more.
\end{tcolorbox}
